\section{Case Study Instrument: Covering Designs}\label{sec:instrument}

\subsection{Problem and verifier asymmetry}

A covering design $C(v,k,t)$ is a collection of $k$-element blocks
from a $v$-element universe such that every $t$-element subset of the
universe is contained in at least one block; the objective is to
minimize the number of blocks. The domain exhibits the
\emph{verifier asymmetry} that guides all our problem choices:
constructing good coverings has occupied specialists for decades, but
verifying a candidate is exact and cheap --- count, over all
$\binom{v}{t}$ $t$-subsets, whether each is covered. Our evaluator
does exactly this count with integer arithmetic; there is no floating
point anywhere in the scoring path.

\subsection{Scoring without leakage}

Feasible solutions with $B$ blocks score
$\mathrm{Sch\ddot{o}nheim}(v,k,t)/B \in (0,1]$, where the
Sch\"onheim bound \cite{schonheim1964} --- a classical,
instance-computable lower bound --- serves as the normalizer. This
choice matters for hygiene: the normalizer injects \emph{no empirical
knowledge} (no archive values, no best-known sizes reach the prompt,
the config, or the scoring function; repository values live in a
reference directory that the loop cannot read). A score of $1.0$
means the Sch\"onheim bound is met, which is rarely possible on large
cells; scores are comparable across cells, enabling multi-instance
fitness sets. Infeasible outputs map into a disjoint band below
every feasible score, graded by the fraction of uncovered
$t$-subsets --- an early timeout that leaves a partial covering still
receives gradient. Format violations (wrong cardinality,
out-of-range or duplicate elements, duplicate blocks) poison the
solution but never crash the evaluator, and work caps guarantee the
instrument terminates on adversarial inputs.

\subsection{Calibrating the instrument against theory and ground truth}

Before any evolution we calibrated the instrument in three layers.
\emph{Curation}: the repository's public data dump was reduced to a
targets table (8{,}759 cells), validated by hard invariants --- the
Sch\"onheim bound never exceeds the recorded lower bound, lower
bounds never exceed recorded sizes, and on the 96 cells with
$k=3, t=2$ the recorded sizes equal the exact Fort--Hedlund value
\cite{fort1958minimal} (96/96). The same pass surfaced 196 rows of
history-hygiene anomalies in auxiliary fields, fixing the project
rule that archive comparisons always target the size field.
\emph{Ground truth}: an exhaustive enumerator (iterative deepening on
the block count with first-uncovered-subset branching and a
$\lceil\text{uncovered}/\binom{k}{t}\rceil$ bound) independently
re-proved nine archive-proven cells, including refuting the
Sch\"onheim value 11 for $C(7,4,3)$ and certifying 12 (450k search
nodes), and proving the gate cell $C(13,3,2)=26$ in milliseconds.
These proven cells act as \emph{gates} in every later experiment: an
evolved program failing to reach a proven optimum signals instrument
or population damage before any claim is risked. \emph{Baseline}:
the human-written seed solver (sampled-candidate greedy, redundancy
elimination, ruin-and-recreate; anytime with atomic output
replacement) reaches the proven optima on the gates but sits far
behind the archive on target cells --- e.g., 1{,}269 vs.\ 620 blocks
on $C(32,8,4)$ at small budgets --- establishing, via the archive
values as reachability certificates, that genuine headroom exists for
evolution rather than a saturated objective.
